\documentclass[11pt]{article}

\usepackage[letterpaper,margin=1in]{geometry}
\usepackage{amsmath,amssymb}
\usepackage{xcolor}
\usepackage{array}
\usepackage[hidelinks]{hyperref}
\usepackage{fancyhdr}

\pagestyle{fancy}
\fancyhead[L]{Team 18105}
\fancyhead[C]{}
\fancyhead[R]{Page \thepage}
\fancyfoot{}
\renewcommand{\headrulewidth}{0.4pt}
\setlength{\parindent}{0pt}
\setlength{\parskip}{8pt}

\begin{document}

% ============================================================================
% PAGE 1 : SUMMARY SHEET
% ============================================================================
\thispagestyle{empty}

% ---- Standard HiMCM summary-sheet three-column header ----
% "B" and "18105" are centered below their titles.
\begin{center}
\begin{minipage}[c]{0.30\textwidth}
\centering
\textbf{Problem Chosen}\\[2pt]
{\Large \textcolor{red}{\textbf{B}}}
\end{minipage}%
\hfill
\begin{minipage}[c]{0.40\textwidth}
\centering
\textbf{2026 HiMCM}\\[2pt]
\textbf{Summary Sheet}
\end{minipage}%
\hfill
\begin{minipage}[c]{0.30\textwidth}
\centering
\textbf{Team Control Number}\\[2pt]
{\Large \textcolor{red}{\textbf{18105}}}
\end{minipage}
\end{center}

\vspace{6pt}
\noindent\hrule
\vspace{12pt}

% ---- Summary (five readable paragraphs) ----
\noindent
On a hot afternoon, commuters waiting at a bus shelter, an elevated platform, or
an open plaza feel the cooled air they paid for slip away into the surrounding
heat. These semi-open spaces are harder to cool than enclosed buildings: with no
sealed boundary, the cooling constantly fights incoming hot air and much of it
blows into empty space. Full air-conditioning is rarely economical or feasible
there, leaving fans and evaporative coolers as the only practical options. A fan
only injects jet velocity and lowers perceived temperature by wind chill, with no
temperature or humidity source; an evaporative cooler adds a cooling sink and a
moisture source, its temperature drop bounded by the wet-bulb depression.

To judge these devices we build a coupled model. A two-dimensional
\textbf{advection--diffusion equation} tracks temperature and humidity under the
wind, and \textbf{moving agents} represent passengers. Two agent modes are needed:
task-locked agents model people who must stay at a fixed waiting position, while
free-waiting agents step locally toward the most comfortable nearby cell. We keep
the coupling one-way, treating agents as field observers, because resolving
two-way plume-obstacle-body interactions would multiply the cost; that refinement
is deferred to the limitations. We convert each agent's perceived temperature
into a two-sided discomfort score that penalizes both overheating and cold
drafts. Summing comfort gained over waiting time and dividing by energy consumed
gives a cooling-efficiency ratio, the first required quantity.

Efficiency alone is not enough: a layout can earn a high ratio with tiny power
yet deliver negligible comfort to a large crowd, or block evacuation routes and
create slip hazards. We therefore add a \textbf{cooling-effectiveness score} built
from six criteria with entropy weighting and \textbf{TOPSIS}; entropy weighting
derives objective weights from sample variance, and we test whether the preferred
layout shifts under priorities such as a cost-focused operator versus a
safety-focused administrator. Layouts that break safety or physical limits are
filtered out first.

Running the full model inside an optimizer is impractical: each high-fidelity
evaluation must solve a transient convection--diffusion field and integrate agent
motion, so hundreds of \textbf{simulated-annealing} iterations would exceed
ordinary laptop runtime. We instead adopt a \textbf{surrogate-assisted search}:
Latin-Hypercube Sampling selects a limited set of high-fidelity runs, these train
a fast random-forest stand-in, simulated annealing searches for the near-optimal
layout on that stand-in alone, and the optimized candidate layout is checked
against the full model.

The general recommendations we report are synthesized by us from cross-comparison
across many cases, not produced directly by the optimizer. We apply the framework
to two venues with contrasting climates: the humid Dhoby Ghaut bus shelter, where
high humidity limits evaporative-cooler performance, and the dry-hot Cartuja Qanat
\'Agora sunken plaza, where \textbf{evaporative cooling} has large theoretical
potential, so the optimal strategy shifts with the environment. The main
limitations are the two-dimensional simplification, the one-way coupling, the
surrogate's prediction error, and the local-optimum nature of simulated
annealing.

\vspace{12pt}
\noindent Keywords: advection-diffusion PDE, agent-based model, surrogate-assisted
simulated annealing, TOPSIS, evaporative cooling, multi-criteria evaluation

% ============================================================================
% PAGE 2 : TABLE OF CONTENTS
% ============================================================================
\clearpage

{\LARGE\bfseries Contents}\par
\vspace{14pt}
\begingroup
\large
\setlength{\parskip}{4pt}

\newcommand{\tocsec}[1]{\par\vspace{10pt}\noindent\textbf{#1}\hfill\par}
\newcommand{\tocsub}[1]{\par\noindent\hspace*{2.2em}#1\dotfill\par}

\tocsec{1\quad Introduction}
\tocsub{1.1\quad Background and Motivation}
\tocsub{1.2\quad Problem Restatement \& Requirements Breakdown}
\tocsub{1.3\quad Overall Framework and Paper Organization}

\tocsec{2\quad Assumptions and Justification}
\tocsub{2.1\quad Physical and Climate Assumptions}
\tocsub{2.2\quad Cooling Device and Energy Assumptions}
\tocsub{2.3\quad Pedestrian and Crowd Assumptions}

\tocsec{3\quad Notation and Key Variables}

\tocsec{4\quad Integrated PDE-Agent Modelling Framework}
\tocsub{4.1\quad Governing Advection-Diffusion PDE}
\tocsub{4.2\quad Dual-Mode Pedestrian Agent Model}
\tocsub{4.3\quad Two-Sided Discomfort Metric \& Cooling Efficiency (Requirement 1)}
\tocsub{4.4\quad Multi-Criteria Cooling Effectiveness with Entropy-TOPSIS (Requirement 4)}

\tocsec{5\quad Surrogate-Assisted Simulated Annealing for Layout Optimization}
\tocsub{5.1\quad Computational Bottleneck of High-Fidelity Simulation}
\tocsub{5.2\quad Latin Hypercube Sampling \& Random Forest Surrogate}
\tocsub{5.3\quad Optimization Search and Back Validation}

\tocsec{6\quad Model Application to Two Distinct Venues (Requirement 3)}
\tocsub{6.1\quad Venue 1: Humid Dhoby Ghaut Bus Shelter}
\tocsub{6.2\quad Venue 2: Dry-Hot Cartuja Qanat \'Agora Plaza}
\tocsub{6.3\quad Cross-Venue Comparison}

\tocsec{7\quad Sensitivity Analysis}

\tocsec{8\quad Strengths, Limitations and Future Work}

\tocsec{9\quad Conclusions}

\tocsec{References}
\tocsec{Appendices}

\endgroup

\end{document}
